%! Change in Linear and Exponential Functions — Unit 2, Lesson 2.2 — Student Packet Cover Sheet
\documentclass[10pt]{article}
\usepackage{precalculus-article}
\usepackage{precalculus-boxes}
\usepackage{ltablex}
\keepXColumns

\begin{document}

% Full-bleed navy banner
\begin{tikzpicture}[remember picture, overlay]
  \fill[navy] (current page.north west)
    rectangle ([yshift=-0.9in]current page.north east);
\end{tikzpicture}
\vspace{-0.8in}

\begin{center}
  {\color{white}\textbf{\LARGE Precalculus}}\\[4pt]
  {\color{skymid}\large\textbf{Unit 2: Exponential and Logarithmic Functions}}\\[6pt]
  {\color{white}\textbf{\large Lesson 2.2 \quad Change in Linear and Exponential Functions}}
\end{center}

\vspace{0.22in}
\namedateperiod
\vspace{0.18in}

\begin{learningtargetbox}
\small
\begin{enumerate}[leftmargin=1.5em, itemsep=5pt]
  \item I can construct a linear function from a slope and a point using
        $f(x) = y_i + m(x - x_i)$, and connect it to an arithmetic sequence.
  \item I can construct an exponential function from a ratio and a point using
        $f(x) = y_i \cdot r^{x - x_i}$, and connect it to a geometric sequence.
  \item I can determine whether a function is linear or exponential by checking whether
        output values over equal-length input intervals have a constant difference or a
        constant ratio.
\end{enumerate}
\end{learningtargetbox}

\vspace{0.14in}

\begin{tocbox}
\small
\renewcommand{\arraystretch}{1.5}
\begin{tabularx}{\linewidth}{@{} c l X r @{}}
\rowcolor{sky}
\textbf{\#} & \textbf{Component} & \textbf{Description} & \textbf{Score} \\
\midrule
1 & Warm-Up        & Spiral review: arithmetic \& geometric sequences from Lesson 2.1 & \blank{1.2cm} \\
2 & Guided Notes   & Linear and exponential functions; point-slope and point-ratio forms & \blank{1.2cm} \\
3 & Group Activity & Build \& classify functions from tables and points (tiered R / A / E) & \blank{1.2cm} \\
4 & Exit Ticket    & Classify from a table; construct from two points; true/false          & \blank{1.2cm} \\
5 & Homework       & Mixed practice: classify, construct, compare, AP-style MC             & \blank{1.2cm} \\
\midrule
  & \multicolumn{2}{r}{\textbf{Total}} & \blank{1.2cm} \\
\end{tabularx}
\end{tocbox}

\end{document}
